% Options for packages loaded elsewhere
\PassOptionsToPackage{unicode}{hyperref}
\PassOptionsToPackage{hyphens}{url}
%
\documentclass[
  11pt,
]{article}
\usepackage{lmodern}
\usepackage{amssymb,amsmath}
\usepackage{ifxetex,ifluatex}
\ifnum 0\ifxetex 1\fi\ifluatex 1\fi=0 % if pdftex
  \usepackage[T1]{fontenc}
  \usepackage[utf8]{inputenc}
  \usepackage{textcomp} % provide euro and other symbols
\else % if luatex or xetex
  \usepackage{unicode-math}
  \defaultfontfeatures{Scale=MatchLowercase}
  \defaultfontfeatures[\rmfamily]{Ligatures=TeX,Scale=1}
  \setmainfont[]{Roboto Light}
\fi
% Use upquote if available, for straight quotes in verbatim environments
\IfFileExists{upquote.sty}{\usepackage{upquote}}{}
\IfFileExists{microtype.sty}{% use microtype if available
  \usepackage[]{microtype}
  \UseMicrotypeSet[protrusion]{basicmath} % disable protrusion for tt fonts
}{}
\makeatletter
\@ifundefined{KOMAClassName}{% if non-KOMA class
  \IfFileExists{parskip.sty}{%
    \usepackage{parskip}
  }{% else
    \setlength{\parindent}{0pt}
    \setlength{\parskip}{6pt plus 2pt minus 1pt}}
}{% if KOMA class
  \KOMAoptions{parskip=half}}
\makeatother
\usepackage{xcolor}
\IfFileExists{xurl.sty}{\usepackage{xurl}}{} % add URL line breaks if available
\IfFileExists{bookmark.sty}{\usepackage{bookmark}}{\usepackage{hyperref}}
\hypersetup{
  pdftitle={Informe del 04 de febrero de 2021},
  hidelinks,
  pdfcreator={LaTeX via pandoc}}
\urlstyle{same} % disable monospaced font for URLs
\usepackage[margin=1in]{geometry}
\usepackage{graphicx,grffile}
\makeatletter
\def\maxwidth{\ifdim\Gin@nat@width>\linewidth\linewidth\else\Gin@nat@width\fi}
\def\maxheight{\ifdim\Gin@nat@height>\textheight\textheight\else\Gin@nat@height\fi}
\makeatother
% Scale images if necessary, so that they will not overflow the page
% margins by default, and it is still possible to overwrite the defaults
% using explicit options in \includegraphics[width, height, ...]{}
\setkeys{Gin}{width=\maxwidth,height=\maxheight,keepaspectratio}
% Set default figure placement to htbp
\makeatletter
\def\fps@figure{htbp}
\makeatother
\setlength{\emergencystretch}{3em} % prevent overfull lines
\providecommand{\tightlist}{%
  \setlength{\itemsep}{0pt}\setlength{\parskip}{0pt}}
\setcounter{secnumdepth}{-\maxdimen} % remove section numbering
\usepackage{titling}
\usepackage{fancyhdr}
\usepackage{lastpage}



\topmargin=-0.45in
\evensidemargin=-1in
\oddsidemargin=0in
\textwidth=6.5in
\textheight=9.0in
\headsep=0.25in

\linespread{1}
 \pagestyle{fancy}
 \fancyhead[LO,LE]{\thetitle}
 \fancyhead[RO,RE]{\includegraphics[width=0.5cm]{./DSLab_logo_2.png}}
 \fancyfoot[LO,LE]{DSLab, methaodos.org, Academia Joven de España \thedate}
 \cfoot{}
 \fancyfoot[RE,RO]{Página\ \thepage\ de\ \protect\pageref*{LastPage}}

\renewcommand\headrulewidth{0.4pt}
\renewcommand\footrulewidth{0.4pt}


\pretitle{%
  \begin{center}
  \LARGE
  \vspace{-2.5cm}
  \includegraphics[width=3.2cm]{./methaodos.png}
  \hspace{3cm}
  \includegraphics[width=3.5cm]{./dslab.png}
  \hspace{3cm}
  \includegraphics[width=3cm]{./aje.png}\\ 
\vspace{1cm}
\textbf{Suficiencia Sanitaria y COVID-19}\\
}
\posttitle{\end{center}}
\usepackage[font=small,skip=8pt]{caption}
\usepackage{float}
\floatplacement{figure}{H}
\renewcommand{\contentsname}{Índice}

\title{Informe del 04 de febrero de 2021}
\author{}
\date{\vspace{-2.5em}}

\begin{document}
\maketitle

\renewcommand{\figurename}{Figura}
\renewcommand{\tablename}{Tabla}

\tableofcontents

\vspace{.5cm}

\begin{center}\rule{0.5\linewidth}{0.5pt}\end{center}

\vspace{.5cm}

Para analizar el comportamiento de la segunda ola de la pandemia
provocada por el virus SARS-CoV-2 y la enfermedad COVID-19 en España, se
han diseñado dos \textbf{Indicadores de Suficiencia Sanitaria (ISS):}
uno acumulado \textbf{(ISSa)} y otro semanal \textbf{(ISSs)}. Además, se
lleva a cabo una estimación del número de \textbf{casos activos} desde
que se inició la pandemia.

\hypertarget{indicador-de-suficiencia-sanitaria-acumulado-issa}{%
\subsection{Indicador de Suficiencia Sanitaria acumulado
(ISSa)}\label{indicador-de-suficiencia-sanitaria-acumulado-issa}}

El ISSa está definido en términos de porcentaje. Desde esa perspectiva,
valores altos del ISSa (cercanos a 100) indicarán que el sistema
sanitario está en disposición de absorber el flujo de diagnósticos
confirmados de coronavirus para que puedan ser tratados según los
protocolos establecidos. Por el contrario, valores bajos del ISSa
(cercanos a 0) sintomatizarán limitaciones que impedirán la correcta
gestión del flujo de casos diagnosticados. Tomando como referencia un
análisis de los valores del ISSa de la primera ola, se deben alcanzar
valores de este indicador por encima del 95\% para poder dar por
controlada la situación de pandemia desde un punto de vista de
suficiencia sanitaria.

Del ISSa se presentan dos estimaciones:

\begin{itemize}
\item
  Una primera estimación de su valor desde el inicio de la segunda ola
  -datada a principios de julio- hasta la actualidad. \textbf{En estos
  momentos, el ISSa de España correspondiente a la segunda ola de la
  pandemia ha descendido hasta el 86.4\%} (ver tabla 1), con una
  tendencia decreciente que se está moderando (ver figura 1, izquierda),
  en lo que puede significar un cambio de tendencia. \textbf{Este
  comportamiento del valor del indicador es coherente con las medidas,
  mayoritariamente reactivas, que se están tomando a lo largo de la
  segunda ola, y especialmente las tomadas después del periodo
  navideño.} Entendemos por \textbf{medidas reactivas} aquellas que se
  toman para dar respuesta a valores extremos de determinados
  indicadores (como por ejemplo la tasa de incidencia a 14 días). El
  denominador común de estas medidas es que se relajan una vez los
  valores de los indicadores dejan de ser extremos.
\item
  Una segunda estimación de su valor a lo largo de toda la pandemia, es
  decir, tomando para su cálculo el periodo completo que abarca desde el
  mes de marzo hasta la actualidad (ver figura 1, derecha). Su valor es
  del 86.25\%. Esta estimación permite ver el avance completo de la
  pandemia.
\end{itemize}

También se calcula el ISSa para cada comunidad autónoma. Se puede
comprobar en la figura 2 que el ISSa de la segunda oleada tiene también
un comportamiento decreciente que se está moderando en la mayor parte de
las autonomías. De hecho, Aragón, Canarias y Cataluña se mantienen en
valores del indicador por encima del 95\%, y otras, como por ejemplo
Extremadura y Murcia, presentan ya un cambio de tendencia. Andalucía
(ISSa = 79.33\%), Comunidad Valenciana (ISSa = 73.88\%) y Galicia (ISSa
= 77.98\%) son las que presentan los valores más bajos del ISSa, aunque
también con una clara moderación de la tendencia decreciente.

\begin{table}[!h]

\caption{\label{tab:tabla}ISS acumulado por CCAA, segunda ola y pandemia.}
\centering
\fontsize{9}{11}\selectfont
\begin{tabular}[t]{l|l|l}
\hline
Comunidad Autónoma & ISSa (\%) 2ª ola & ISSa (\%) pandemia\\
\hline
Andalucia & 79.33 & 80.23\\
\hline
Aragón & > 95 & > 95\\
\hline
Asturias & 87.32 & 88.35\\
\hline
Baleares & 91.91 & 92.38\\
\hline
C. Valenciana & 73.88 & 75.18\\
\hline
Canarias & > 95 & > 95\\
\hline
Cantabria & 92.57 & 93.63\\
\hline
Castilla La Mancha & 83.72 & 86.09\\
\hline
Castilla y León & 83.45 & 85.75\\
\hline
Cataluña & > 95 & > 95\\
\hline
Ceuta & > 95 & > 95\\
\hline
Extremadura & 91.15 & 91.8\\
\hline
Galicia & 77.98 & 80.54\\
\hline
La Rioja & 83.4 & 86.57\\
\hline
Madrid & 89.99 & 91.39\\
\hline
Melilla & > 95 & > 95\\
\hline
Murcia & 87.11 & 87.54\\
\hline
Navarra & > 95 & > 95\\
\hline
País Vasco & 92.38 & 93.63\\
\hline
España & 86.4 & 86.25\\
\hline
\end{tabular}
\end{table}

\begin{figure}
\centering
\includegraphics{20210204_COVID19_DSLAB_informe_ejecutivo_files/figure-latex/fig_iss3-1.pdf}
\caption{\label{fig:fig_iss3} ISS acumulado, segunda ola y pandemia.}
\end{figure}

\vspace{0.2cm}

\begin{figure}
\centering
\includegraphics{20210204_COVID19_DSLAB_informe_ejecutivo_files/figure-latex/fig17a_res-1.pdf}
\caption{\label{fig:fig17a_res} ISS acumulado en las Comunidades
Autónomas, segunda ola.}
\end{figure}

\vspace{0.2cm}

\begin{figure}
\centering
\includegraphics{20210204_COVID19_DSLAB_informe_ejecutivo_files/figure-latex/fig17a_res2-1.pdf}
\caption{\label{fig:fig17a_res} ISS acumulado en las Comunidades
Autónomas, desde el inicio de la pandemia.}
\end{figure}

\newpage
\setcounter{page}{6}

\hypertarget{indicador-de-suficiencia-sanitaria-semanal-isss}{%
\subsection{Indicador de Suficiencia Sanitaria semanal
(ISSs)}\label{indicador-de-suficiencia-sanitaria-semanal-isss}}

El ISSs es un indicador que toma valores positivos, siendo clave superar
un valor de 1 para garantizar la suficiencia del sistema. Sucesivos
valores semanales por debajo de 1 implican posibles saturaciones del
sistema sanitario, mientras que si el indicador se sitúa por encima de 1
evidencia mejores niveles de suficiencia. Para que el sistema sanitario
esté en condiciones de estabilizarse es necesario que el ISSa mantenga
una tendencia creciente y que el ISSs se sitúe de manera sostenida por
encima de 1. \textbf{Desde una perspectiva nacional, el ISSs de la
segunda ola ha vuelto a bajar del valor de 1 (ISSs = 0.78, ver tabla 2),
manteniendo una tendencia oscilante en torno a 1 desde el mes de
septiembre (ver figura 4).} Desagregando por comunidades autónomas,
todas ellas presentan un comportamiento similar (ver figura 5). En el
momento este indicador se estabilice por encima de 1, el número de casos
resueltos semanalmente superará al de casos diagnosticados. \textbf{Los
valores oscilantes del ISSs, tanto a nivel nacional como autonómico, son
coherentes con las oscilaciones del indicador acumulado, el ISSa.}

\begin{table}[!h]

\caption{\label{tab:tabla2}ISS semanal por CCAA, segunda ola y pandemia.}
\centering
\fontsize{9}{11}\selectfont
\begin{tabular}[t]{l|r|r}
\hline
Comunidad Autónoma & ISSs 2ª ola & ISSs pandemia\\
\hline
Andalucia & 0.61 & 0.61\\
\hline
Aragón & 1.03 & 1.03\\
\hline
Asturias & 0.29 & 0.29\\
\hline
Baleares & 0.45 & 0.45\\
\hline
C. Valenciana & 0.78 & 0.78\\
\hline
Canarias & 0.89 & 0.89\\
\hline
Cantabria & 0.78 & 0.78\\
\hline
Castilla La Mancha & 0.97 & 0.95\\
\hline
Castilla y León & 0.64 & 0.64\\
\hline
Cataluña & 1.29 & 1.25\\
\hline
Ceuta & 0.93 & 0.93\\
\hline
Extremadura & 1.75 & 1.75\\
\hline
Galicia & 0.51 & 0.51\\
\hline
La Rioja & 0.70 & 0.70\\
\hline
Madrid & 0.64 & 0.64\\
\hline
Melilla & 0.84 & 0.84\\
\hline
Murcia & 1.43 & 1.43\\
\hline
Navarra & 0.90 & 0.90\\
\hline
País Vasco & 0.31 & 0.31\\
\hline
España & 0.78 & 0.78\\
\hline
\end{tabular}
\end{table}

\begin{figure}
\centering
\includegraphics{20210204_COVID19_DSLAB_informe_ejecutivo_files/figure-latex/fig_iss2-1.pdf}
\caption{\label{fig:fig_iss2} ISS semanal, segunda ola y pandemia.}
\end{figure}

\begin{figure}
\centering
\includegraphics{20210204_COVID19_DSLAB_informe_ejecutivo_files/figure-latex/figura_ISSSD_ccaa-1.pdf}
\caption{\label{fig:figura_ISSSD_ccaa} ISS semanal por Comunidades
Autónomas.}
\end{figure}

\clearpage

\setcounter{page}{9}

\hypertarget{estimaciuxf3n-de-casos-activos}{%
\subsection{Estimación de casos
activos}\label{estimaciuxf3n-de-casos-activos}}

Las \textbf{figuras 6 y 7} muestran tendencias que resumen de forma muy
ajustada la evolución de la pandemia en España. En concreto, presentan
\textbf{la estimación del número de casos diagnosticados que permanecen
activos cada semana}, tanto a nivel nacional como por comunidades
autónomas. \textbf{La mayoría las comunidades autónomas muestran
escenarios crecientes después del periodo navideño en lo relativo al
número de casos activos. Está tendencia se está moderando e incluso
invirtiendo en algunas comunidades autónomas, como por ejemplo Aragón,
Cataluña, Canarias, Extremadura y Murcia. Esta inflexión en la tendencia
creciente, de confirmarse en el tiempo, unida a la aplicación de medidas
preventivas (como por ejemplo la vacunación), deberían implicar una
mejora de la suficiencia sanitaria y, en consecuencia, del control de la
pandemia.}

\vspace{0.2cm}

\begin{figure}
\centering
\includegraphics{20210204_COVID19_DSLAB_informe_ejecutivo_files/figure-latex/fig1-1.pdf}
\caption{\label{fig:fig17a_res} Estimación de casos activos en España
desde el inicio de la pandemia.}
\end{figure}

\begin{figure}
\centering
\includegraphics{20210204_COVID19_DSLAB_informe_ejecutivo_files/figure-latex/figura_activos_ccaa-1.pdf}
\caption{\label{fig:figura_activos_ccaa} Estimación de casos activos por
Comunidades Autónomas.}
\end{figure}

\newpage
\setcounter{page}{11}

\hypertarget{definiciones}{%
\subsection{Definiciones}\label{definiciones}}

Definiciones básicas:

\begin{itemize}
\item
  \(Confirmados_s\): número de diagnósticos confirmados acumulados a
  fecha \(s\).
\item
  \(NuevosCasos_{s} = Confirmados_{s} - Confirmados_{s-1}\): nuevos
  contagiados la semana \(s\).
\item
  \(Fallecidos_s\): número de fallecimientos acumulados a fecha \(s\).
\item
  \(FallecidosSemanales_s = Fallecidos_{s} - Fallecidos_{s-1}\): número
  de fallecimientos la semana \(s\).
\item
  \(Recuperados_s\): número de casos recuperados acumulados a fecha
  \(s\).
\item
  \(RecuperadosSemanales_{s} = Recuperados_{s} - Recuperados_{s-1}\):
  número de recuperados en la semana \(s\).
\item
  \(ActivosSemanales_s = Confirmados_s - Fallecidos_s - Recuperados_s\).
  Casos activos al final de la semana \(s\).
\item
  \(AtendidosSemanales_{s} = RecuperadosSemanales_s + FallecidosSemanales_s\):
  Pacientes atendidos la semana \(s\).
\item
  \(Atendidos_{s} = Recuperados_s + Fallecidos_s\): Pacientes atendidos
  acumulados a fecha \(s\).
\item
  \textbf{ISS acumulado a semana \(s\)}
  \[ISS\_Acumulado_s = \frac{Atendidos_s}{Confirmados_s} \times 100\]
  Este indicador mide el estrés acumulado del sistema. Cuanto más
  próximo a 100, más desestresado está el sistema.
\item
  \textbf{ISS Semanal, semana \(s\)}
\end{itemize}

\[
ISS_s = \frac{AtendidosSemanales_s}{NuevosCasos_{s}}
\]

\vspace{4cm}

\begin{center}
\includegraphics{cc.png}
\end{center}

\vspace{0.2cm}

\begin{center}
Esta obra está bajo una licencia de Creative Commons 
\end{center}

\begin{center}
Reconocimiento-CompartirIgual 4.0 Internacional.
\end{center}

\end{document}
